% 分野別類題: 定数係数2階線形非同次微分方程式（定数変化法） 解答例
% コンパイル: TEXINPUTS="../../../00_common:$TEXINPUTS" latexmk -xelatex answers.tex
\documentclass[a4paper,11pt]{article}
\usepackage{preamble}
\usepackage{shoakubox}

\begin{document}

\kamokumei{類題演習 解答例 --- 定数係数2階線形非同次微分方程式（定数変化法）}

\daimon{【解法】定数変化法（ロンスキアン利用）の解き方と導出}

\noindent
2階線形非同次微分方程式
\[
y'' + p(x)\,y' + q(x)\,y = f(x)
\]
に対応する同次方程式 $y''+p(x)y'+q(x)y=0$ の1次独立な基本解を $y_{1}(x), y_{2}(x)$ とする。未定係数法は $f(x)$ が多項式・指数関数・三角関数の組合せの場合にしか使えないため、$f(x) = \tan x$ や $1/x$ のような一般の関数に対しては\textbf{定数変化法}を用いる。

同次方程式の一般解は $y = C_{1}y_{1}+C_{2}y_{2}$ であるが、定数変化法では $C_{1}, C_{2}$ を $x$ の関数 $u_{1}(x), u_{2}(x)$ に置き換え、
\[
y = u_{1}(x)\,y_{1}(x) + u_{2}(x)\,y_{2}(x)
\]
の形の特殊解を探す。$u_{1}, u_{2}$ には自由度があるので、まず
\[
u_{1}'\,y_{1} + u_{2}'\,y_{2} = 0 \tag{$*$}
\]
という補助条件を課す。このとき
\[
y' = u_{1}y_{1}' + u_{2}y_{2}'
\]
となる（$u_{1}', u_{2}'$ を含む項が $(*)$ により消える）。さらに微分すると
\[
y'' = u_{1}'y_{1}' + u_{1}y_{1}'' + u_{2}'y_{2}' + u_{2}y_{2}''
\]
これらを元の方程式に代入すると、$y_{1}, y_{2}$ が同次方程式の解であることから $u_{1}(y_{1}''+py_{1}'+qy_{1}) + u_{2}(y_{2}''+py_{2}'+qy_{2}) = 0$ の部分は消え、
\[
u_{1}'\,y_{1}' + u_{2}'\,y_{2}' = f(x) \tag{$**$}
\]
が残る。$(*)$ と $(**)$ を $u_{1}', u_{2}'$ に関する連立1次方程式とみて解くと、係数行列の行列式はロンスキアン
\[
W(x) = \begin{vmatrix} y_{1} & y_{2} \\ y_{1}' & y_{2}' \end{vmatrix} = y_{1}y_{2}' - y_{2}y_{1}'
\]
であり（$y_{1}, y_{2}$ が1次独立なので $W \neq 0$）、クラメルの公式より
\[
u_{1}' = -\frac{y_{2}\,f(x)}{W(x)}, \qquad u_{2}' = \frac{y_{1}\,f(x)}{W(x)}
\]
を得る。これらを積分して $u_{1}(x), u_{2}(x)$ を求め、特殊解 $y_{p} = u_{1}y_{1}+u_{2}y_{2}$ を得る。一般解は
\[
y = C_{1}y_{1} + C_{2}y_{2} + y_{p}
\]
である。初期条件が与えられた場合は、これに代入して $C_{1}, C_{2}$ を定める。

\clearpage
\begin{mondaiboxS}{問1}
$y'' + y = \tan x \quad \left(-\dfrac{\pi}{2} < x < \dfrac{\pi}{2}\right)$ の一般解を求めよ。
\end{mondaiboxS}
\kaitou
同次方程式 $y''+y=0$ の基本解は $y_{1}=\cos x,\ y_{2}=\sin x$。ロンスキアンは
\[
W = \cos x\cdot\cos x - \sin x\cdot(-\sin x) = \cos^{2}x+\sin^{2}x = 1
\]
よって
\[
u_{1}' = -y_{2}\tan x = -\sin x\tan x = -\frac{\sin^{2}x}{\cos x} = -\sec x + \cos x,
\qquad
u_{2}' = y_{1}\tan x = \cos x\tan x = \sin x
\]
積分して
\[
u_{1} = -\ln|\sec x+\tan x| + \sin x, \qquad u_{2} = -\cos x
\]
（積分定数は $0$ ととる。）特殊解は
\[
y_{p} = u_{1}y_{1}+u_{2}y_{2}
= \left(\sin x - \ln|\sec x+\tan x|\right)\cos x + (-\cos x)\sin x
= -\cos x\,\ln|\sec x+\tan x|
\]
したがって一般解は
\[
\boxed{\; y = C_{1}\cos x + C_{2}\sin x - \cos x\,\ln|\sec x+\tan x| \;}
\qquad (C_{1}, C_{2}\text{ は任意定数})
\]
（検算: $h(x)=\ln|\sec x+\tan x|$ とおくと $h'(x)=\sec x$ であるから、$y_{p}=-\cos x\,h(x)$ より $y_{p}'=\sin x\,h-1$、$y_{p}''=\cos x\,h+\tan x$。よって $y_{p}''+y_{p}=\tan x$ が成り立つ。）

\clearpage
\begin{mondaiboxS}{問2}
$y'' + y = \dfrac{1}{\cos x} \quad \left(-\dfrac{\pi}{2} < x < \dfrac{\pi}{2}\right)$ の一般解を求めよ。
\end{mondaiboxS}
\kaitou
同次方程式の基本解・ロンスキアンは問1と同じく $y_{1}=\cos x,\ y_{2}=\sin x,\ W=1$。$f(x)=\sec x$ より
\[
u_{1}' = -y_{2}\sec x = -\sin x\sec x = -\tan x,
\qquad
u_{2}' = y_{1}\sec x = \cos x\sec x = 1
\]
積分して
\[
u_{1} = \ln|\cos x| = \ln(\cos x), \qquad u_{2} = x
\]
（この区間では $\cos x>0$。）特殊解は
\[
y_{p} = u_{1}y_{1}+u_{2}y_{2} = \cos x\,\ln(\cos x) + x\sin x
\]
したがって一般解は
\[
\boxed{\; y = C_{1}\cos x + C_{2}\sin x + \cos x\,\ln(\cos x) + x\sin x \;}
\qquad (C_{1}, C_{2}\text{ は任意定数})
\]
（検算: $y_{p}'=-\sin x\ln(\cos x)+x\cos x$、$y_{p}''=-\cos x\ln(\cos x)+\dfrac{\sin^{2}x}{\cos x}+\cos x-x\sin x$ より
$y_{p}''+y_{p}=\dfrac{\sin^{2}x+\cos^{2}x}{\cos x}=\dfrac{1}{\cos x}$ が成り立つ。）

\clearpage
\begin{mondaiboxS}{問3}
$y'' - 2y' + y = \dfrac{e^{x}}{x}, \quad y(1)=e,\ y'(1)=e \quad (x>0)$ を解け。
\end{mondaiboxS}
\kaitou
特性方程式 $(r-1)^{2}=0$ より重解 $r=1$ で、同次方程式の基本解は $y_{1}=e^{x},\ y_{2}=xe^{x}$。ロンスキアンは
\[
W = y_{1}y_{2}' - y_{2}y_{1}' = e^{x}(e^{x}+xe^{x}) - xe^{x}\cdot e^{x} = e^{2x}
\]
$f(x)=\dfrac{e^{x}}{x}$ より
\[
u_{1}' = -\frac{y_{2}f}{W} = -\frac{xe^{x}\cdot \frac{e^{x}}{x}}{e^{2x}} = -1,
\qquad
u_{2}' = \frac{y_{1}f}{W} = \frac{e^{x}\cdot\frac{e^{x}}{x}}{e^{2x}} = \frac{1}{x}
\]
積分して $u_{1} = -x,\ u_{2} = \ln x$（$x>0$）。特殊解は
\[
y_{p} = u_{1}y_{1}+u_{2}y_{2} = -xe^{x} + xe^{x}\ln x
\]
ここで $-xe^{x}$ は同次解 $y_{2}=xe^{x}$ の定数倍なので任意定数に吸収でき、一般解は
\[
y = Ae^{x} + Bxe^{x} + xe^{x}\ln x = e^{x}\left(A + Bx + x\ln x\right)
\]
とおける。初期条件を適用する。$y(1) = e(A+B) = e$ より $A+B=1$。また
\[
y' = e^{x}(A+Bx+x\ln x) + e^{x}(B+\ln x + 1)
\]
なので、$x=1$ を代入すると $y'(1)=e\left[(A+B)+(B+1)\right]=e(A+2B+1)=e$ より $A+2B=0$。連立して解くと $A+B=1,\ A+2B=0$ より $B=-1,\ A=2$。したがって
\[
\boxed{\; y = e^{x}\left(2 - x + x\ln x\right) \;} \qquad (x>0)
\]
（検算: $g(x)=2-x+x\ln x$ とおくと $g'(x)=\ln x$。$y=e^{x}g$ より $y'=e^{x}(g+\ln x)$、$y''=e^{x}\left(g+2\ln x+\dfrac{1}{x}\right)$。よって
\[
y''-2y'+y = e^{x}\left[\left(g+2\ln x+\tfrac{1}{x}\right)-2(g+\ln x)+g\right] = e^{x}\cdot\frac{1}{x} = \frac{e^{x}}{x}
\]
が成り立つ。また $y(1)=e(2-1+0)=e$、$y'(1)=e(g(1)+\ln 1)=e(1+0)=e$ となり初期条件も満たす。）

\end{document}
